%=======================================================================
%  New Structural and Arithmetic Constraints on Finite-Support
%  Periodic Highways of Langton's Ant
%
%  Author: Atharva Jillhewar
%
%  SELF-CONTAINED: paste this single file into Overleaf and press
%  "Recompile" (pdfLaTeX).  The bibliography is embedded, so there is NO
%  separate .bib file and NO BibTeX step required.  References and
%  appendices render in one compile (Overleaf runs the 2 pdfLaTeX passes
%  needed to resolve \cite and \ref automatically).
%
%  >>> Before posting: fill in the \address and \email placeholders below.
%=======================================================================
\documentclass[11pt,reqno]{amsart}

\usepackage[utf8]{inputenc}
\usepackage[T1]{fontenc}
\usepackage{amsmath,amssymb,amsthm,mathtools}
\usepackage{enumitem}
\usepackage{booktabs}
\usepackage{xcolor}
\usepackage{listings}
\usepackage[margin=1in]{geometry}
\usepackage[colorlinks=true,linkcolor=blue!55!black,citecolor=teal!60!black,
            urlcolor=purple!55!black]{hyperref}

%------------------------------ theorem environments
\theoremstyle{plain}
\newtheorem{theorem}{Theorem}[section]
\newtheorem{lemma}[theorem]{Lemma}
\newtheorem{proposition}[theorem]{Proposition}
\newtheorem{corollary}[theorem]{Corollary}
\theoremstyle{definition}
\newtheorem{definition}[theorem]{Definition}
\newtheorem{remark}[theorem]{Remark}
\newtheorem{example}[theorem]{Example}
\newtheorem{conjecture}[theorem]{Conjecture}
\newtheorem{problem}[theorem]{Problem}

%------------------------------ macros
\newcommand{\Z}{\mathbb{Z}}
\newcommand{\F}{\mathbb{F}}
\newcommand{\che}{\chi}
\newcommand{\drift}{d}
\newcommand{\growth}{g}
\newcommand{\turnR}{\mathrm{R}}
\newcommand{\turnL}{\mathrm{L}}
\newcommand{\vE}{\mathsf{E}}
\DeclareMathOperator{\rank}{rank}
\DeclareMathOperator{\supp}{supp}

%------------------------------ code listing style (ASCII-faithful Lean 4)
\lstdefinestyle{lean}{
  basicstyle=\ttfamily\footnotesize,
  keywordstyle=\bfseries,
  commentstyle=\itshape\color{gray},
  morekeywords={def,theorem,lemma,by,match,with,structure,namespace,end,
    import,fun,let,in,Bool,List,true,false,rw,simp,induction,cases,exact,
    refine,intro,rfl},
  showstringspaces=false,
  columns=fullflexible,
  breaklines=true,
  frame=single,
  framesep=5pt,
  rulecolor=\color{gray!40},
  backgroundcolor=\color{gray!5},
}

\title[Constraints on Langton-ant highways]{New Structural and Arithmetic Constraints on\\ Finite-Support Periodic Highways of Langton's Ant}

\author{Atharva Jillhewar}
% >>> EDIT the next two lines with your affiliation and email before posting:
\address{(affiliation)}
\email{(email)}

\date{\today}

\keywords{Langton's ant, turmite, cellular automata, highway conjecture,
symbolic dynamics, reversible dynamics, medial graph, Tait graph, interlacement}
\subjclass[2020]{Primary 37B15; Secondary 37B10, 05C10, 68Q80, 68R15}

\begin{document}

\begin{abstract}
For the classical two-colour Langton's ant on $\Z^2$ with finite initial black
support, we establish the first structural and arithmetic constraints on the
possible \emph{periodic highways}. Our principal results are the following.
\emph{(1) An even-winding theorem}: no finite-support pattern can drift to an
exact translate of itself with restored heading and zero net change in the number
of black cells; consequently every finite-support periodic highway must build a
wake whose net black growth is a strictly positive multiple of four. \emph{(2) A
signed mod-four wake-residue identity}: the checkerboard-signed strand bases of any
periodic highway sum to $2\pmod 4$ in each residue class of every nonzero drift
coordinate, whence the sharp lower bound $\growth\ge 2\max(|a|,|b|)$ tying growth
$\growth$ to drift $(a,b)$. \emph{(3) A Tait-graph conjugacy}: the ant is exactly
conjugate to edge-toggling on a medial (Tait) graph, with a cycle-rank surgery
identity $E+\Delta C=2\beta$. \emph{(4) A two-endpoint collision-chain parity
theorem} and a \emph{touch-graph interlacement rank formula} for the periodic
trace. \emph{(5) A two-engine, rank-audited exhaustive computation} excluding every
finite-support periodic highway of nonzero drift and period at most $48$: the
classical ant has no short-period highway other than its period-$104$ one. We also
give a decidable periodic-realisability criterion with an explicit finite seed,
strengthening the forbidden-word program for the ant's trace subshift. The
algebraic core of the even-winding and residue theorems is machine-checked in Lean
$4$. Every result is a necessary condition on a finite-support periodic highway or
an exact finite computation; none resolves the (open) convergence conjecture, and
we delimit precisely what a solution must still supply.
\end{abstract}

\maketitle

\tableofcontents

%=======================================================================
\section{Introduction}
%=======================================================================

\emph{Langton's ant} is the turmite on the square lattice $\Z^2$ defined as
follows. Each cell is coloured white or black; a pointer, the \emph{ant}, occupies
a cell and faces one of the four cardinal directions. At each step the ant reads
the colour of its current cell: on white it turns $90^\circ$ right, on black it
turns $90^\circ$ left; it then flips the colour of the current cell and moves one
unit forward in its new heading. Starting from the all-white plane the ant produces,
after a transient of $9977$ steps, a self-similar diagonal structure---the
\emph{highway}---that repeats with least period $104$ and translates the ant by
$(\pm2,\pm2)$ each period.

Throughout we impose the \emph{finite-support} condition: only finitely many cells
are initially black. The empirical universality of the highway is the subject of a
well-known open problem.

\begin{conjecture}[Highway conjecture]\label{conj:main}
For every finite initial black set and every initial pose, the orbit of Langton's
ant eventually enters a translate, rotation, reflection, and phase shift of the
standard period-$104$ highway.
\end{conjecture}

\subsection{Prior art}
The rigorous facts about the classical ant are few. Every trajectory is
\emph{unbounded} (Bunimovich--Troubetzkoy \cite{BunimovichTroubetzkoy1992}; the
underlying checkerboard/parity structure is already in Gale--Propp--Sutherland--Troubetzkoy
\cite{GaleProppSutnerTroubetzkoy1995}). Prediction problems for the ant are
computationally universal and undecidable, so no finite time cutoff can settle
Conjecture~\ref{conj:main} \cite{GajardoMoreiraGoles2002}. Gajardo
\cite{Gajardo2003symbolic} gives an algebraic characterisation of the forbidden
factors of the ant's trace subshift. Recent progress concerns \emph{generalised}
ants: some rules admit infinitely many highways \cite{GajardoLutfallaRao2024},
others have genuinely non-highway emergent behaviour \cite{Lutfalla2025sideways},
and quantitative confinement bounds are available \cite{Etse2026confined}. None of
this literature constrains what a \emph{classical} finite-support periodic highway
can look like: there is no published net-growth congruence, drift-versus-growth
arithmetic, translator-exclusion theorem, or medial-graph analysis for the classical
ant. The present paper supplies exactly that missing structural layer.

\subsection{Results}
A \emph{repeating trace} is an infinite turn word $w^\omega$ that is realised as the
trace of some finite-support colouring (Section~\ref{sec:criterion}); its
\emph{drift} $\drift=(a,b)$ is the per-period displacement and its \emph{growth}
$\growth$ the per-period change in the number of black cells. Our contributions,
with an explicit account of their novelty:

\begin{enumerate}[leftmargin=2.4em,label=\textbf{(\arabic*)}]
\item \textbf{Even-winding theorem} (Theorem~\ref{thm:evenwinding}): no
finite-support clean translator with nonzero drift exists. Hence every
finite-support periodic highway has growth $\growth\in4\Z_{>0}$
(Corollary~\ref{cor:growth4}). Bunimovich--Troubetzkoy excludes only \emph{bounded}
repeatable orbits; a zero-growth glider is unbounded and is not covered by their
theorem, so the strict positivity is new. (The divisibility by $4$ alone is the
elementary parity Lemma~\ref{lem:invariant}.)

\item \textbf{Signed mod-four wake-residue identity} (Theorem~\ref{thm:residue}) and
the sharp bound $\growth\ge2\max(|a|,|b|)$ (Corollary~\ref{cor:P16}). We are not
aware of any exact per-residue arithmetic identity, or any sharp linear drift--growth
bound, in the literature.

\item \textbf{Tait-graph conjugacy} (Theorem~\ref{thm:conjugacy}) and the cycle-rank
surgery identity $E+\Delta C=2\beta$ (Theorem~\ref{thm:cyclerank}). No previous
Langton-ant work uses medial/Tait graphs or loop merge/split surgery; the only
precedent is the far weaker checkerboard partition of
\cite{GaleProppSutnerTroubetzkoy1995}.

\item \textbf{Collision-chain parity} (Theorem~\ref{thm:collision}) and the
\emph{touch-graph interlacement rank} formula (Theorem~\ref{thm:touchgraph}), the
latter built on Traldi's interlacement/transition machinery
\cite{Traldi2016binary,Traldi2012transition}. (We write ``touch-graph
interlacement'' to avoid a clash with the unrelated ``interlaced word'' of
\cite{Lutfalla2025sideways}.)

\item \textbf{Small-period exclusion} (Theorem~\ref{thm:periodexcl}): a two-engine,
rank-audited exhaustive search excludes every finite-support periodic highway of
nonzero drift and period at most $48$. Existing computation enumerates highways that
\emph{occur}; proving non-existence below a period bound is new.
\end{enumerate}

Section~\ref{sec:criterion} also gives a decidable periodic-realisability criterion
with an explicit finite seed, strengthening the forbidden-word program of
\cite{Gajardo2003symbolic} from local factor-admissibility to global periodic
finite-support realisability; only that strengthening is claimed as new. The
elementary heading/black-count invariant (Lemma~\ref{lem:invariant}) is folklore,
implicit in \cite{BunimovichTroubetzkoy1992,GaleProppSutnerTroubetzkoy1995}.

\subsection{Scope}
Every theorem below is a \emph{necessary condition} on a finite-support periodic
highway, or an exact finite computation. None asserts that any orbit converges, and
therefore none resolves Conjecture~\ref{conj:main}. What we prove is that any
finite-support periodic highway must build a positive wake divisible by four, must
satisfy the residue identities and the density bound $\growth\ge2\|\drift\|_\infty$,
and cannot have period at most $48$ unless it is the standard highway. The single
geometric ingredient separating these results from a full solution is stated
precisely in Section~\ref{sec:gap}.

%=======================================================================
\section{Setup and elementary facts}\label{sec:setup}
%=======================================================================

A \emph{state} is a triple $s=(B,p,q)$ with $B\subset\Z^2$ finite (the black set),
$p\in\Z^2$ (the ant position) and $q\in\Z/4\Z$ the heading, numbered $0=N$, $1=E$,
$2=S$, $3=W$, with unit vectors
\[
u_0=(0,1),\quad u_1=(1,0),\quad u_2=(0,-1),\quad u_3=(-1,0).
\]
The \emph{turn} at $s$ is $\turnR$ (right) with $\varepsilon=+1$ if $p\notin B$, and
$\turnL$ (left) with $\varepsilon=-1$ if $p\in B$. One step of the dynamics is
\begin{equation}\label{eq:step}
s=(B,p,q)\ \longmapsto\ \bigl(B\,\triangle\,\{p\},\ p+u_{q+\varepsilon},\ q+\varepsilon\bigr),
\end{equation}
where $\triangle$ denotes symmetric difference and indices are read modulo $4$. The
\emph{trace} of an orbit is the infinite word of turns encountered. Since
\eqref{eq:step} is invertible (step backward one cell, recolour, undo the turn), the
dynamics is reversible and orbits never merge. Write $b_t=|B_t|$ and let
$\che(x,y)=(-1)^{x+y}$ be the checkerboard sign.

\begin{proposition}[Checkerboard/HV partition {\cite{GaleProppSutnerTroubetzkoy1995}}]
\label{prop:hv}
The quantity $\bigl(q_t-(p_{t,x}+p_{t,y})\bigr)\bmod 2$ is constant along every
orbit. Consequently each cell is entered by the ant always horizontally or always
vertically. Calling a cell \emph{vertical} (resp.\ \emph{horizontal}) accordingly,
at a vertical cell the two vertical incident edges are the only admissible arrival
edges and the two horizontal ones the only departure edges, and vice versa; every
admissible directed lattice edge thus carries a fixed orientation.
\end{proposition}

\begin{proof}
Each step changes the heading by an odd quarter-turn, so $q_t\bmod2$ alternates;
the displacement $u_{q_{t+1}}$ has odd $\ell^1$-length, so $(p_{t,x}+p_{t,y})\bmod2$
alternates as well. Their difference is therefore invariant. Fixing a cell $z$: at
every visit the checkerboard parity of $z$ is the same, so $q_t\bmod2$ is the same,
i.e.\ the arrival axis is fixed. The edge orientations follow because an arrival
edge of one endpoint is a departure edge of the other.
\end{proof}

\begin{proposition}[Alternating visits]\label{prop:alt}
At any fixed cell the encountered turns alternate strictly between $\turnR$ and
$\turnL$. All but finitely many cells begin with $\turnR$.
\end{proposition}

\begin{proof}
Each visit flips the cell's colour and the turn is a function of the colour, so the
turn alternates. A cell white at its first visit begins with $\turnR$; only the
finitely many initially black cells can begin with $\turnL$.
\end{proof}

\begin{theorem}[Unboundedness {\cite{BunimovichTroubetzkoy1992}}]\label{thm:unbounded}
Every finite-support orbit visits infinitely many cells.
\end{theorem}

\begin{proof}[Proof sketch]
Were the visited set finite, reversibility would make the orbit periodic in the full
state. Choose a visited cell $z$ maximising $x+y$; its north and east neighbours are
unvisited. If $z$ is horizontal it is entered from the west, and its two successive
(white then black) visits leave south and then north, forcing a visit to an unvisited
corner neighbour---a contradiction; the vertical case is the $90^\circ$ rotation.
\end{proof}

\begin{lemma}[Heading/black-count invariant; folklore]\label{lem:invariant}
For every $t$,
\begin{equation}\label{eq:invariant}
q_t-b_t\equiv q_0-b_0\pmod 4,\qquad p_{t+1}-p_t=u_{\,q_0-b_0+b_{t+1}} .
\end{equation}
A heading-resetting period changes $b$ by a multiple of $4$.
\end{lemma}

\begin{proof}
A white/$\turnR$ step increments both $q$ and $b$ by $1$; a black/$\turnL$ step
decrements both. Hence $q_t-b_t$ is constant. Since the move follows the turn,
$q_{t+1}\equiv q_0-b_0+b_{t+1}$, giving the displacement formula. If $q_N=q_0$ then
$b_N-b_0\equiv 0\pmod4$.
\end{proof}

%=======================================================================
\section{A constructive periodic-realisability criterion}\label{sec:criterion}
%=======================================================================

We first record the exact test that turns ``is $w^\omega$ a highway?'' into a finite
computation. Fix an initial heading $q_0$ and a proposed turn word
$w=w_0\cdots w_{N-1}\in\{\turnR,\turnL\}^N$. Compute the induced path
$p_0,\dots,p_N$ from $w$ (Proposition~\ref{prop:hv} makes each step well defined).
Suppose $q_N=q_0$ and $\drift=p_N-p_0\neq0$. Define an equivalence on phases by
\[
i\sim j\iff p_i-p_j\in\Z\drift,
\]
and in each class $I$ write $p_i=b_I+a_i\drift$ for a fixed representative $b_I$ and
integers $a_i$. For $n\ge\min_I a_i$, let $S_{I,n}$ be the subword of all $w_i$
($i\in I$) with $a_i\le n$, listed in increasing order of the key $(n-a_i,\,i)$.

\begin{theorem}[Periodic-realisability criterion]\label{thm:criterion}
The infinite word $w^\omega$ is the trace of Langton's ant from some finite-support
colouring if and only if, for every class $I$:
\begin{enumerate}[label=\textup{(\roman*)}]
\item every $S_{I,n}$ with $\min_I a_i\le n\le\max_I a_i$ alternates strictly between
$\turnR$ and $\turnL$; and
\item the stabilised word $S_{I,\max_I a_i}$ begins with $\turnR$.
\end{enumerate}
When these hold, an explicit finite black seed producing $w^\omega$ is
\begin{equation}\label{eq:seed}
B=\bigl\{\,b_I+n\drift\ :\ n<\textstyle\max_I a_i\ \text{and}\ S_{I,n}\ \text{begins with}\ \turnL\,\bigr\}.
\end{equation}
\end{theorem}

\begin{proof}
The occurrence of phase $i\in I$ in period $k\ge0$ sits at
$p_i+k\drift=b_I+(a_i+k)\drift$, which equals $b_I+n\drift$ exactly when $k=n-a_i\ge0$;
chronological order among these occurrences is the lexicographic order of $(n-a_i,i)$.
Thus $S_{I,n}$ is precisely the chronological turn sequence at the physical cell
$b_I+n\drift$ under $w^\omega$.

\emph{Necessity.} Each visit flips the cell, so its turn sequence alternates,
proving (i). For $n\ge\max_I a_i$ all phases of $I$ are present and the ordering no
longer depends on $n$: the sequence has stabilised, and the infinitely many cells
$b_I+n\drift$ ($n\ge\max_I a_i$) share this first turn. By Proposition~\ref{prop:alt}
only finitely many cells begin with $\turnL$, so the stable first turn is $\turnR$,
proving (ii).

\emph{Sufficiency.} Colour exactly the set \eqref{eq:seed} black. At each boundary
cell $b_I+n\drift$ with $n<\max_I a_i$ the chosen colour reproduces the first symbol
of $S_{I,n}$, and strict alternation supplies every subsequent symbol because the
cell flips at each visit. At every stabilised cell ($n\ge\max_I a_i$) condition (ii)
makes the first encountered turn $\turnR$, matching its white initial colour.
Induction on chronological time shows the ant follows $w^\omega$ forever; since
$q_N=q_0$ and the position advances by $\drift\neq0$ each period, this is a highway.
\end{proof}

\begin{remark}\label{rem:onesort}
There is an equivalent single-sequence test: within a class, sort all entries by
decreasing $a_i$, breaking ties by increasing $i$; every $S_{I,n}$ is a suffix of the
result obtained by deleting whole high-level blocks. It therefore suffices to check
that this one stabilised word alternates and begins with $\turnR$, at the cost of one
sort per class. A class and its level are located in constant time by reducing one
coordinate modulo the corresponding coordinate of $\drift$. This is the checker used
in all computations of Section~\ref{sec:computation}; it strengthens the
forbidden-word characterisation of the ant's trace subshift
\cite{Gajardo2003symbolic} to a decidable test for global periodic finite-support
realisability with a constructive seed. A worked instance (the standard highway) is
given in Appendix~\ref{app:example}.
\end{remark}

%=======================================================================
\section{The Tait-graph conjugacy}\label{sec:tait}
%=======================================================================

We recast the dynamics on a medial (Tait) graph, exposing topology invisible in the
raw colouring. Normalise the HV partition of Proposition~\ref{prop:hv} so that
even cells receive vertical arrivals. In the \emph{frozen} routing (the permutation
of directed states obtained by ignoring colour flips) the admissible directed arcs
form clockwise directed four-cycles around one checkerboard class of unit plaquettes.
Label a plaquette by its lower-left corner and take these plaquettes as the vertices
of a rotated square lattice. Each ant cell $z=(x,y)$ lies at the crossing of exactly
two white four-cycles and corresponds to the \emph{Tait edge}
\begin{equation}\label{eq:tait}
e_z=\begin{cases}\{(x,y-1),(x-1,y)\}, & x+y\ \text{even},\\[2pt]
\{(x-1,y-1),(x,y)\}, & x+y\ \text{odd}.\end{cases}
\end{equation}
Let $H_B=\{e_z:z\in B\}$ be the finite edge subgraph induced by the black cells.

\begin{theorem}[Boundary-loop conjugacy]\label{thm:conjugacy}
White and black at $z$ are the two smoothings of the corresponding directed crossing.
Consequently:
\begin{enumerate}[label=\textup{(\alph*)}]
\item the frozen routing loops are exactly the boundary components of a regular
neighbourhood of $H_B$, together with the untouched four-cycles;
\item the ant traverses its current boundary loop and then toggles the edge $e_z$; and
\item each step either merges two frozen loops into one or splits one into two.
\end{enumerate}
This is an exact conjugacy, not a heuristic.
\end{theorem}

Let $E=|B|$ be the number of occupied Tait edges, $V$ the number of incident Tait
vertices, $k$ the number of nonempty components, and $\beta=E-V+k$ the cycle rank.

\begin{theorem}[Cycle-rank surgery identity]\label{thm:cyclerank}
A regular neighbourhood of $H_B$ has $k+\beta=E-V+2k$ boundary components. Writing
$\Delta C$ for the change in loop count relative to the $V$ trivial plaquette loops,
\begin{equation}\label{eq:cyclerank}
\Delta C=E-2V+2k,\qquad E+\Delta C=2\beta .
\end{equation}
Adding a bridge merges two loops and fixes $\beta$; adding a cycle edge splits a loop
and raises $\beta$ by one; removing a bridge splits a loop and fixes $\beta$;
removing a non-bridge merges two loops and lowers $\beta$ by one.
\end{theorem}

\begin{proof}
A regular neighbourhood of a graph with $V$ vertices, $E$ edges, and $k$ components is
a compact surface with $k$ pieces and Euler characteristic $V-E$; its boundary
consists of $k+\beta$ circles, where $\beta=E-V+k$ is the first Betti number.
Formula~\eqref{eq:cyclerank} is the difference from the $V$ trivial plaquette loops.
The four surgery cases are the standard effect on the boundary of adding or deleting
a bridge (Betti number unchanged; boundary count changes by $\mp1$) versus a cycle
edge (Betti number changes by $\pm1$; opposite merge/split), read off from
Theorem~\ref{thm:conjugacy}(c).
\end{proof}

\begin{proposition}[Advanced records are bridges]\label{prop:record}
Once a bounding-box record lies outside the finite initial support, its first visit
is white, and a new east/north/west/south record turns south/east/north/west. In the
Tait graph the corresponding edge has a previously untouched endpoint, hence is a
bridge. Therefore sufficiently advanced record creation never raises $\beta$; any
unbounded production of cycle rank occurs strictly inside already-explored territory.
\end{proposition}

\begin{remark}[No monotone invariant]\label{rem:nomonotone}
The active-loop length, its signed area, $E$, $V$, $k$, and $\beta$ all fail to be
monotone, with explicit counterexamples on the blank orbit (e.g.\ the active loop
drops from length $20$ to $4$ at step $4$; the cycle rank takes every value in
$\{1,\dots,5\}$ during one mature period). Compact states realise both signs of every
independent local change of $(E,V,k,\beta)$, so no nonconstant affine combination of
these four quantities is universally monotone. This rules out the most direct
potential-function proofs of Conjecture~\ref{conj:main}.
\end{remark}

%=======================================================================
\section{Additive charges and clean translators}\label{sec:charges}
%=======================================================================

\begin{definition}\label{def:translator}
A \emph{clean finite translator with nonzero drift} is a finite state for which, for
some $N>0$ and $\drift\in\Z^2\setminus\{0\}$,
\[
B_N=B_0+\drift,\qquad p_N=p_0+\drift,\qquad q_N=q_0 .
\]
The word \emph{clean} means the entire black set translates: no stationary residue
and no newly emitted wake. This is not the ordinary highway, which grows a wake.
\end{definition}

\subsection{Laurent charges and axis alignment}
Work in the Laurent polynomial ring $\F_2[X^{\pm1},Y^{\pm1}]$, an integral domain.
Let $V_t=1$ if $q_t$ is vertical and $H_t=1-V_t$, and set
\[
A_t=\sum_{(x,y)\in B_t}X^x+V_t\,X^{x_t},\qquad
D_t=\sum_{(x,y)\in B_t}Y^y+H_t\,Y^{y_t}.
\]

\begin{proposition}\label{prop:laurent}
$A_t$ and $D_t$ are conserved. Hence a clean translator with $\drift=(a,b)$ satisfies
$A=X^aA$ and $D=Y^bD$; since $A(1)=|B|+V$ and $D(1)=|B|+H$ cannot both vanish in
$\F_2$ (because $V+H=1$), the drift coordinates $a$ and $b$ are not both nonzero.
Every clean translator is axis-aligned.
\end{proposition}

\begin{proof}
Conservation of $A_t$: on a step from a vertical cell, the toggle at column $x_t$
cancels the disappearing vertical correction $X^{x_t}$; on a step from a horizontal
cell, the post-turn motion is vertical, so $x_{t+1}=x_t$ and the toggle cancels the
newly appearing correction. Translation covariance then gives $A=X^aA$. In the
integral domain $\F_2[X^{\pm1},Y^{\pm1}]$, $a\neq0$ forces $A=0$ and $b\neq0$ forces
$D=0$; but $A(1)+D(1)=2|B|+1=1\neq0$.
\end{proof}

\begin{proposition}[Additive-charge cocycle]\label{prop:cocycle}
Let $m\ge2$. A weight $w:\Z^2\to\Z/m\Z$ admits heading corrections
$F_q:\Z^2\to\Z/m\Z$ making $\sum_{z\in B}w(z)+F_q(p)$ conserved if and only if $mw=0$
and
\begin{equation}\label{eq:cocycle}
w(x{+}1,y{+}1)+w(x{+}1,y)+w(x,y{+}1)+w(x,y)=0\quad(\bmod\ m).
\end{equation}
The step relations are then, for the conserved quantity,
\begin{equation}\label{eq:pot}
F_{q+1}(z+u_{q+1})-F_q(z)=-w(z)\ \ (\turnR),\qquad
F_{q-1}(z+u_{q-1})-F_q(z)=+w(z)\ \ (\turnL).
\end{equation}
Over $\F_2$ the solutions of \eqref{eq:cocycle} are $w(x,y)=\alpha(x)+\beta(y)$.
\end{proposition}

\begin{proof}
On an $\turnR$ step the black set gains $z$ (weight $+w(z)$) and the state moves
$(z,q)\mapsto(z+u_{q+1},q+1)$; conservation gives the first relation in
\eqref{eq:pot}, and the $\turnL$ case is symmetric. Eliminating the four $F_q$ from
\eqref{eq:pot} around a unit plaquette yields exactly \eqref{eq:cocycle}, whose
$\F_2$ solution space is the separable weights $\alpha(x)+\beta(y)$.
\end{proof}

A mod-four refinement of Proposition~\ref{prop:laurent} (using the quadratic solution
$Q(x,y)=x^2+2xy$ of \eqref{eq:cocycle} over $\Z/4$) strengthens the drift of a clean
translator to $(4a,0)$ or $(0,4a)$, with period divisible by $8$ and $N\ge8|\drift|$;
we use only axis alignment below and refer to \cite{Repo} for the refinement.

\subsection{The strand decomposition}
Let $\growth=\#\turnR-\#\turnL$ over one period. Lemma~\ref{lem:invariant} gives
$\growth\equiv0\pmod4$, and $\growth\ge0$ because $\growth$ equals the number of
\emph{odd} stabilised translation classes of Theorem~\ref{thm:criterion} (a class is
odd iff its stabilised word has odd length, contributing one net black cell per
period). We record the endpoint structure used in Section~\ref{sec:residue}.

\begin{proposition}[Canonical endpoint decomposition]\label{prop:decomp}
The one-period signed toggle pattern $D:\Z^2\to\Z$ (with $D(z)=\#\{\turnR\text{ at }z\}-\#\{\turnL\text{ at }z\}$ over one period) decomposes as
\begin{equation}\label{eq:decomp}
D=S+(T_\drift-1)A,\qquad S=\sum_{i=1}^{\growth}\delta_{c_i},
\end{equation}
where $(T_\drift A)(z)=A(z-\drift)$, $A$ is a finitely supported \emph{moving widget},
and $c_1,\dots,c_\growth\in\Z^2$ are the bases of the $\growth$ growing strands, one
per odd stabilised class. For the standard word this is a $2$-cell stationary residue,
an $11$-cell widget, and $12$ strands, giving $|B_n|=13+12n$ at period boundaries.
\end{proposition}

If $\growth=0$ then $S=0$ and, after the finite stationary residue separates from the
path, deleting it leaves a clean translator; Sections~\ref{sec:evenwinding}
and~\ref{sec:residue} prove this is impossible, by two independent routes.

%=======================================================================
\section{The even-winding theorem}\label{sec:evenwinding}
%=======================================================================

This section proves that no clean translator exists, via a cylinder homology
argument. The one step that was only sketched in earlier accounts---the local
parity identity relating the turn count at a cell to its horizontal edge
parity---is proved in full as Lemma~\ref{lem:localparity}.

\begin{theorem}[Even-winding theorem]\label{thm:evenwinding}
No finite-support clean translator with nonzero drift exists.
\end{theorem}

For the remainder of the section fix a hypothetical clean translator. By
Proposition~\ref{prop:laurent} and the mod-four refinement its drift is axis-aligned
with a coordinate divisible by four; after a rotation and reflection write
$\drift=(L,0)$ with $L\in4\Z_{>0}$. Reversibility and translation covariance give,
for all $n\in\Z$ and phases $j$,
\[
A_{nN+j}=A_j+n\drift,\qquad p_{nN+j}=p_j+n\drift .
\]
Quotient the bi-infinite trace by translation through $\drift$ to obtain a closed
walk on the cylinder $C_L=(\Z/L\Z)\times\Z$; because $L$ is even the HV type of
Proposition~\ref{prop:hv} descends to $C_L$. For an undirected cylinder edge, its
\emph{multiplicity} is the number of traversals by the projected one-period walk.

\begin{lemma}[Visit count]\label{lem:visit}
Every vertex $v$ of $C_L$ is visited $n_v=2k_v$ times in one quotient period, with
exactly $k_v$ turns $\turnR$ and $k_v$ turns $\turnL$.
\end{lemma}

\begin{proof}
Let $z$ be a physical representative of $v$. Each phase $j$ with $p_j\in v+\Z\drift$
has a unique level $m_j$ with $p_j=z+m_j\drift$, and phase $j$ in period $-m_j$ visits
$z$; conversely every visit to $z$ over the bi-infinite orbit arises once this way.
Hence the visits to $v$ in one quotient period biject with the visits to $z$ over the
translated bi-infinite orbit. At each fixed phase the black pattern is a finite
translate of $A_0$, so $z$ is white for all sufficiently early and all sufficiently
late times. Its turn sequence therefore alternates, begins with $\turnR$, ends with
$\turnL$, and contains equally many of each.
\end{proof}

At a vertex $v$ let $H_{v,-},H_{v,+}$ be the multiplicities of the west and east
horizontal edges and $V_{v,-},V_{v,+}$ those of the south and north vertical edges.
By Proposition~\ref{prop:hv} one pair carries all arrivals and the other all
departures, and each visit is one arrival plus one departure, so
\begin{equation}\label{eq:balance}
H_{v,-}+H_{v,+}=V_{v,-}+V_{v,+}=2k_v .
\end{equation}

\begin{lemma}[Edge parities]\label{lem:edgeparity}
The horizontal-edge multiplicity has a constant parity $h_y\in\F_2$ along each row
$y$. Every vertical-edge multiplicity is even; write $V_e=2G_e$ with $G_e\in\Z_{\ge0}$.
\end{lemma}

\begin{proof}
Reducing the horizontal half of \eqref{eq:balance} modulo $2$ gives
$H_{v,-}\equiv H_{v,+}$, so consecutive horizontal edges in a row share parity;
call it $h_y$. The vertical half gives $V_{v,-}\equiv V_{v,+}$, so vertical parity is
constant up each column. The projected walk has finite vertical support, so a
vertical edge far enough below or above it has multiplicity $0$; constancy then forces
every vertical multiplicity even.
\end{proof}

\begin{lemma}[Local parity]\label{lem:localparity}
At every vertex $v=(x,y)$ of $C_L$, $\;k_v\equiv h_y\pmod2$.
\end{lemma}

\begin{proof}
Suppose $v$ is vertical (the horizontal case is the $90^\circ$ rotation). By
Proposition~\ref{prop:hv} the ant arrives at $v$ along a vertical edge---from the
south moving north, or from the north moving south---and departs along a horizontal
edge. Record each visit by (arrival direction, turn). Using
$u_{q+\varepsilon}$ from \eqref{eq:step}: an arrival from the south ($\turnR$) departs
east, from the south ($\turnL$) departs west, from the north ($\turnR$) departs west,
from the north ($\turnL$) departs east. Let $a_{\turnR},a_{\turnL}$ count arrivals
from the south with the two turns and $b_{\turnR},b_{\turnL}$ those from the north.
Because the east edge is only ever traversed \emph{out} of $v$ (its other endpoint is
horizontal, hence cannot depart along it), its multiplicity equals the number of east
departures:
\[
H_{v,+}=a_{\turnR}+b_{\turnL}.
\]
The north edge is only ever traversed \emph{into} $v$, with multiplicity
$V_{v,+}=b_{\turnR}+b_{\turnL}$. The total turn count is $k_v=a_{\turnR}+b_{\turnR}$.
Hence
\[
H_{v,+}-k_v=(a_{\turnR}+b_{\turnL})-(a_{\turnR}+b_{\turnR})=b_{\turnL}-b_{\turnR}
\equiv b_{\turnL}+b_{\turnR}=V_{v,+}\pmod2 .
\]
By Lemma~\ref{lem:edgeparity}, $H_{v,+}$ has parity $h_y$ and $V_{v,+}$ is even; therefore
$h_y\equiv k_v\pmod2$.
\end{proof}

\begin{proof}[Proof of Theorem~\ref{thm:evenwinding}]
Halving the vertical half of \eqref{eq:balance} and using $V_e=2G_e$ gives the exact
identity $k_{x,y}=G_{x,y-1/2}+G_{x,y+1/2}$, where $G_{x,y\pm1/2}$ denote the halved
multiplicities of the vertical edges below and above $v$. Fix a column $x$ and sum
Lemma~\ref{lem:localparity} over all rows:
\begin{equation}\label{eq:telescope}
\sum_y h_y\ \equiv\ \sum_y k_{x,y}\ =\ \sum_y\bigl(G_{x,y-1/2}+G_{x,y+1/2}\bigr)\ \equiv\ 0\pmod2,
\end{equation}
because each interior half-edge $G$ appears exactly twice and the two exterior terms
vanish by finite vertical support.

Now cut $C_L$ along a vertical seam between two adjacent columns. The algebraic
winding number $w$ of the projected closed walk is the signed number of seam
crossings; there is exactly one horizontal seam edge in each row $y$, of multiplicity
$\equiv h_y\pmod2$, and signs vanish modulo $2$, so
\[
w\equiv\sum_y h_y\equiv0\pmod2
\]
by \eqref{eq:telescope}. On the other hand the lift of the closed walk runs from $p$
to $p+(L,0)$, so it winds exactly once around $C_L$: $w=1$. This contradiction proves
the theorem for drift $(L,0)$; reflection and the $90^\circ$ rotation cover the
remaining axis-aligned drifts.
\end{proof}

\begin{corollary}[Growth is a positive multiple of four]\label{cor:growth4}
Every finite-support eventually-translating periodic trajectory has net black growth
\begin{equation}\label{eq:growth4}
\growth\in4\Z_{>0}.
\end{equation}
In particular every finite-support periodic highway builds a nonempty wake.
\end{corollary}

\begin{proof}
Zero drift confines the ant to the finitely many positions of one period, excluded by
Theorem~\ref{thm:unbounded}. Negative growth cannot recur forever since $b_t\ge0$. A
heading-resetting period has $\growth\equiv0\pmod4$ (Lemma~\ref{lem:invariant}). By
Proposition~\ref{prop:decomp}, a $\growth=0$ trace of nonzero drift separates into a
finite stationary residue and a clean moving packet; deleting the residue yields a
clean translator, contradicting Theorem~\ref{thm:evenwinding}. Hence $\growth>0$, and
$\growth\in4\Z_{>0}$.
\end{proof}

The standard highway has $\growth=12$, consistent with \eqref{eq:growth4}. The
telescope \eqref{eq:telescope} is the finite algebraic core machine-checked in Lean
(Appendix~\ref{app:lean}).

%=======================================================================
\section{The signed mod-four wake-residue theorem}\label{sec:residue}
%=======================================================================

We now give a second, independent proof that no zero-growth periodic highway exists,
which additionally pins the arithmetic relation between growth and drift. Recall the
strand bases $c_1,\dots,c_\growth$ of Proposition~\ref{prop:decomp}.

\begin{theorem}[Signed mod-four wake-residue identity]\label{thm:residue}
Let a finite-support repeating trace have restored heading, nonzero drift
$\drift=(a,b)$, growth $\growth$, and decomposition \eqref{eq:decomp}. Then:
\begin{enumerate}[label=\textup{(\roman*)}]
\item if $a\neq0$, then for every residue $r\in\Z/|a|\Z$,
\begin{equation}\label{eq:res-a}
\sum_{i:\;c_{i,x}\equiv r\ (\mathrm{mod}\ |a|)}\che(c_i)\ \equiv\ 2\pmod4;
\end{equation}
\item if $b\neq0$, the analogous sum over each residue of $c_{i,y}$ modulo $|b|$ is
$\equiv2\pmod4$;
\item if $a=0$, the signed sum $\sum_{i:\,c_{i,x}=x_0}\che(c_i)$ over each exact
column $x_0$ is $\equiv0\pmod4$; symmetrically for exact rows when $b=0$.
\end{enumerate}
\end{theorem}

The proof rests on a monodromy computation for the potentials of
Proposition~\ref{prop:cocycle}.

\begin{lemma}[Diagonal increment and vertical periodicity]\label{lem:diag}
Fix $a\neq0$ and any $\alpha:\Z/|a|\Z\to\Z/4\Z$, and set $w(x,y)=\che(x,y)\,\alpha(x)$.
Then $w$ satisfies $4w=0$ and \eqref{eq:cocycle}, so potentials $F_q$ as in
\eqref{eq:pot} exist, and
\begin{equation}\label{eq:diag}
F_N(x{+}1,y{+}1)-F_N(x,y)=-\che(x,y)\bigl(\alpha(x)+\alpha(x{+}1)\bigr),
\end{equation}
while $F_N$ is $2$-periodic in $y$.
\end{lemma}

\begin{proof}
Since $\che(x{+}1,y{+}1)=\che(x,y)$ and $\che(x{+}1,y)=\che(x,y{+}1)=-\che(x,y)$, the
four-corner sum \eqref{eq:cocycle} telescopes to
$\che(x,y)[\alpha(x{+}1)-\alpha(x{+}1)-\alpha(x)+\alpha(x)]=0$, and $4w=0$ is clear;
so $F_q$ exist. Compose the $\turnR$ step at $(x,y)$ heading $N$ (giving
$F_E(x{+}1,y)-F_N(x,y)=-w(x,y)$) with the $\turnL$ step at $(x{+}1,y)$ heading $E$
(giving $F_N(x{+}1,y{+}1)-F_E(x{+}1,y)=+w(x{+}1,y)$). Adding,
\[
F_N(x{+}1,y{+}1)-F_N(x,y)=-w(x,y)+w(x{+}1,y)=-\che(x,y)\alpha(x)-\che(x,y)\alpha(x{+}1),
\]
which is \eqref{eq:diag}. The relation $F_S=F_N-2w$ (from composing the two vertical
steps) gives the same increment along $(1,-1)$; combining the $(1,1)$ and $(1,-1)$
increments over a closed square shows $F_N(x,y{+}2)=F_N(x,y)$.
\end{proof}

\begin{proof}[Proof of Theorem~\ref{thm:residue}]
Fix $a\neq0$, $\alpha$, and $w$ as in Lemma~\ref{lem:diag}. Restored heading gives
$\growth\equiv0\pmod4$, so the period is even and $a+b$ is even; hence translation by
$\drift$ preserves $\che$, and $w(z+\drift)=w(z)$. Telescoping \eqref{eq:diag} across
one full horizontal residue cycle (length dividing $|a|$ after using the $2$-periodicity
in $y$) yields the monodromy
\begin{equation}\label{eq:monodromy}
F_q(z+\drift)-F_q(z)\ \equiv\ 2\sum_{r\in\Z/|a|\Z}\alpha(r)\pmod4,
\end{equation}
for every heading $q$ and either sign of $a$; the telescoped total is
$-2\che(z)\sum_r\alpha(r)$ and $-2\equiv2\pmod4$.

Charge conservation of $\sum_{z\in B}w(z)+F_q(p)$ over one period gives
\begin{equation}\label{eq:conservation}
\sum_{z}D(z)\,w(z)+F_q(p+\drift)-F_q(p)=0,
\end{equation}
where $D$ is the signed toggle pattern of Proposition~\ref{prop:decomp}. Because $w$
is $\drift$-periodic and $(T_\drift A)(z)=A(z-\drift)$, the widget terms cancel:
$\sum_z(T_\drift A)(z)w(z)=\sum_zA(z)w(z)$. Substituting $D=S+(T_\drift-1)A$ and the
monodromy \eqref{eq:monodromy} into \eqref{eq:conservation},
\begin{equation}\label{eq:reduced}
\sum_{i=1}^{\growth}\che(c_i)\,\alpha(c_{i,x})+2\sum_{r\in\Z/|a|\Z}\alpha(r)\ \equiv\ 0\pmod4 .
\end{equation}
Taking $\alpha=\mathbf 1_{\{r\}}$ turns \eqref{eq:reduced} into
$\sum_{i:\,c_{i,x}\equiv r}\che(c_i)+2\equiv0$, i.e.\ \eqref{eq:res-a}, using
$-2\equiv2$. Rotating by $90^\circ$ gives (ii). For $a=0$, take a finitely supported
exact-column $\alpha:\Z\to\Z/4\Z$; then $b$ is even, $w$ is $\drift$-periodic, the
vertical $2$-periodicity makes the monodromy vanish, and exact-column indicators give
(iii).
\end{proof}

\begin{corollary}[Strand-density bound]\label{cor:P16}
Every finite-support repeating trace with drift $\drift=(a,b)$ has
\begin{equation}\label{eq:P16}
\growth\ \ge\ 2\max(|a|,|b|).
\end{equation}
In particular $\growth=0$ is impossible---an independent proof of
Theorem~\ref{thm:evenwinding}. If a residue fiber contains exactly two strand bases,
they share checkerboard parity; for $\drift=(\pm2,\pm2)$ and $\growth=4$ the
$2\times2$ residue-count matrix is $\left[\begin{smallmatrix}2&0\\0&2\end{smallmatrix}\right]$
or its transpose.
\end{corollary}

\begin{proof}
A sum of $m$ signs $\pm1$ congruent to $2\pmod4$ is nonzero, with $m$ positive and
even. Hence each of the $|a|$ residue classes modulo $|a|$ contains at least two
strand bases, giving $\growth\ge2|a|$; symmetrically $\growth\ge2|b|$. If $\growth=0$
some class is empty while $\drift\neq0$, contradicting \eqref{eq:res-a}. A two-element
fiber has $\che$-sum $\pm2$, forcing equal signs.
\end{proof}

Unlike the even-winding argument, Theorem~\ref{thm:residue} uses neither axis
alignment nor the clean-packet reduction; it applies to every repeating trace and
additionally fixes the drift--growth arithmetic. The identities \eqref{eq:res-a} and
the bound \eqref{eq:P16} were verified independently on all $172{,}092$ heading-reset
nonzero-drift words through length $18$ and on the standard highway and its dihedral
images (Appendix~\ref{app:repro}).

%=======================================================================
\section{Collision chains and touch-graph interlacement}\label{sec:interlace}
%=======================================================================

We encode a finite cell set as edges of the bipartite row--column incidence graph:
the cell $(x,y)$ is the edge $C_xR_y$ between a column vertex $C_x$ and a row vertex
$R_y$. To a directed state $s=(p,q)$ assign the token $\eta(s)=C_{p_x}$ if $q$ is
vertical and $\eta(s)=R_{p_y}$ if $q$ is horizontal.

\begin{theorem}[Two-endpoint collision-chain parity]\label{thm:collision}
Over $\F_2$, the boundary (set of odd-degree vertices) of the row--column incidence
graph of the symmetric difference $B_s\triangle B_t$ equals $\eta(s)+\eta(t)$.
Intermediate tokens telescope across any collision checkpoints; consequently at most
one defect component is non-Eulerian, and a same-pose return ($\eta(s)=\eta(t)$) has
only Eulerian defect components. Every completed $\turnR$-to-$\turnL$ black lifetime
has an Eulerian odd-visit component through its paired cell and therefore involves at
least four distinct cells.
\end{theorem}

\begin{proof}[Proof sketch]
One step toggles the single edge $C_{p_x}R_{p_y}$, changing the degree parity of
exactly $C_{p_x}$ and $R_{p_y}$; the pre- and post-step tokens $\eta$ differ by the
one incidence-vertex whose axis matches the heading, so the boundary contribution of
a single step is $\eta(s)+\eta(s')$. Summing over a run telescopes to the two
endpoints. A non-Eulerian component needs a nonempty boundary, so at most one exists,
and it is absent at a same-pose return. The four-cell lower bound follows because an
Eulerian graph with a nonzero cycle through the paired cell has at least four edges of
even degree at each incidence vertex.
\end{proof}

For the periodic branch, pair every stabilised $\turnR$ with its following $\turnL$
into a single $4$-valent temporal vertex; the wake $\turnR$ events (those with no
matching $\turnL$ in a period) become degree-two $H/V$ switches. Suppressing the
switches yields a connected $4$-regular Euler system, to which Traldi's modified
interlacement theory applies \cite{Traldi2016binary,Traldi2012transition}.

\begin{theorem}[Touch-graph interlacement rank]\label{thm:touchgraph}
Let $m$ be the number of paired lifetimes, $c$ the number of closed $H/V$ smoothing
circuits, and $\kappa$ the number of cycles in the $2$-regular wake-edge subgraph
$W$. The labelled boundary touch graph $T$ satisfies $\dim Z(T)=m-c+1$ and its
wake-contraction $T/W$ satisfies $\dim Z(T/W)=m-c-\kappa+1$, fitting the exact
sequence
\[
0\longrightarrow Z(W)\longrightarrow Z(T)\longrightarrow Z(T/W)\longrightarrow0.
\]
The lifetime-interval rows span $Z(T/W)$ but have codimension at most $\kappa$ in
$Z(T)$. On the standard trace $m=46$, $c=5$, $\kappa=1$, projected rank $41$, full
boundary rank $42$.
\end{theorem}

Thus ordinary circuit-nullity closes the contracted graph but cannot supply the
missing $\kappa$-dimensional wake-cycle restriction. At fixed $\growth=4$ there is an
explicit abstract family with $c\to\infty$ and unbounded contracted rank, so no bound
depending only on $\growth$ exists; controlling that kernel with Langton-specific
voltage data is an open step (Section~\ref{sec:gap}).

%=======================================================================
\section{Small-period exclusion}\label{sec:computation}
%=======================================================================

All searches are exact: integer coordinates, exact finite sets, and no probabilistic
equality tests. Two independently implemented engines are used---a reference
implementation of the criterion of Theorem~\ref{thm:criterion}, and a compiled
depth-first search that additionally prunes with Corollary~\ref{cor:growth4} and the
residue identity of Theorem~\ref{thm:residue}. Both engines positively accept the
standard word. The methodology, per-shard records, and cross-audit are described in
Appendix~\ref{app:repro}.

\begin{theorem}[No short-period highway other than $104$]\label{thm:periodexcl}
No finite-support periodic highway of nonzero drift has period at most $48$.
\end{theorem}

\begin{proof}[Verification]
Odd periods restore no heading (an odd number of quarter turns is never
$0\bmod4$), and Corollary~\ref{cor:growth4} excludes zero growth at every period, so
it suffices to search even periods for positive-growth traces. The evidence is:
\begin{enumerate}[label=\textup{(\alph*)}]
\item a complete enumeration of the legal trace tree through period $32$
($46{,}185{,}421$ feasible nodes over $32$ prefix shards; zero highways);
\item exact positive-growth searches at periods $34,36,38,40$, visiting respectively
$9{,}726{,}176$, $19{,}781{,}050$, $41{,}972{,}884$, and $201{,}924{,}597$ nodes,
with zero highways;
\item two-engine, rank-by-rank audited searches at periods $42,44,46,48$. At the
largest, period $48$, both engines visited the identical $8{,}677{,}026{,}370$ nodes
and $451{,}962{,}870$ leaves over all $2^{16}$ prefix ranks with no node cap; the
residue-pruned engine rejected the vast majority of leaves via
Theorem~\ref{thm:residue} and applied the exact criterion to the survivors; both
returned zero certificates, and all per-rank node, leaf, and prune counters agree.
(At period $46$ the two engines likewise agreed on $3{,}463{,}441{,}745$ nodes and
$98{,}568{,}824$ leaves with zero certificates.)
\end{enumerate}
Table~\ref{tab:excl} summarises. Together these exclude every finite-support periodic
highway of nonzero drift and period at most $48$.
\end{proof}

\begin{table}[h]
\centering
\begin{tabular}{@{}llrr@{}}
\toprule
Period & Method & Nodes & Highways \\
\midrule
$\le 32$ & complete legal enumeration & $46{,}185{,}421$ & $0$\\
$34$ & positive-growth & $9{,}726{,}176$ & $0$\\
$36$ & positive-growth & $19{,}781{,}050$ & $0$\\
$38$ & positive-growth & $41{,}972{,}884$ & $0$\\
$40$ & positive-growth & $201{,}924{,}597$ & $0$\\
$42,44$ & two-engine, rank-audited & (see records) & $0$\\
$46$ & two-engine, rank-audited & $3{,}463{,}441{,}745$ & $0$\\
$48$ & two-engine, rank-audited & $8{,}677{,}026{,}370$ & $0$\\
\bottomrule
\end{tabular}
\caption{Exhaustive exclusion of finite-support periodic highways of nonzero drift.
Each entry is complete over its parameter range with no node cap; the period-$46$ and
$48$ rows are cross-validated by two independent engines whose per-rank counters agree.}
\label{tab:excl}
\end{table}

\begin{proposition}[Hamming isolation]\label{prop:hamming}
Every period-$104$ word distinct from the standard one and within Hamming distance
four of it fails the criterion: all $5{,}356$ distance-two and $4{,}598{,}126$
distance-four mutations are excluded. Odd Hamming distance cannot restore the heading.
\end{proposition}

By \cite{GajardoMoreiraGoles2002}, no finite period bound can settle
Conjecture~\ref{conj:main}; Theorem~\ref{thm:periodexcl} is therefore a lower bound on
the minimal period of any nonstandard highway, not a proof of the conjecture.

%=======================================================================
\section{The remaining gap}\label{sec:gap}
%=======================================================================

The results above constrain the arithmetic and topology of any finite-support
periodic highway, but do not force convergence. Two honest obstructions delimit the
gap.

\begin{problem}[Bounded retreat / bounded active core]\label{prob:retreat}
Show that for every finite-support initial configuration there is a time after which
the ant stays within a bounded corridor behind its record frontier; equivalently in
the Tait picture, that the active boundary component at or ahead of the current record
has uniformly bounded size and that components shed behind an advancing frontier are
never reabsorbed.
\end{problem}

Under either formulation the frontier window has finitely many states modulo
translation, determinism forces eventual periodicity, and Theorem~\ref{thm:criterion}
classifies the resulting cycle. No proof is known, and the following show why the easy
routes fail: unboundedness (Theorem~\ref{thm:unbounded}) permits an orbit that keeps
enlarging a two-dimensional region; no scalar Tait invariant is monotone
(Remark~\ref{rem:nomonotone}); and a local pumping argument is invalid because equal
local cores may connect through arbitrarily distant explored territory, turning a
later merge into a split (Theorem~\ref{thm:conjugacy}).

\begin{problem}[Periodic classification]\label{prob:periodic}
Show that every positive-growth periodic trace admitted by finite support is the
standard period-$104$ trace up to symmetry and phase.
\end{problem}

At fixed growth $\growth=4$ the residue skeletons of Theorem~\ref{thm:residue} and the
touch-graph rank of Theorem~\ref{thm:touchgraph} leave unbounded free parameters, and
explicit heading-reset countermodels satisfy every incidence and residue condition yet
fail the cross-period criterion. A resolution therefore needs a single-tour
chronology/interlacement theorem valid for unbounded period---genuinely more than the
incidence data used here. A \emph{disproof} of Conjecture~\ref{conj:main} would
instead exhibit an infinite certificate: a nonstandard periodic word passing
Theorem~\ref{thm:criterion}, or a finite self-enlarging Tait gadget whose state recurs
with a parameter increased. No such certificate has been found.

%=======================================================================
\section{Concluding remarks}\label{sec:conclusion}
%=======================================================================

We have given the first structural and arithmetic theory of the classical
finite-support periodic highway: a decidable realisability criterion, an exact
medial-graph conjugacy with surgery calculus, two independent proofs that no
zero-growth clean translator exists (so every highway grows a wake of size
$\in4\Z_{>0}$), the sharp strand-density bound $\growth\ge2\|\drift\|_\infty$, a
two-endpoint collision-chain parity theorem and a touch-graph rank formula, and an
exhaustive exclusion of all periodic highways of period at most $48$. The algebraic
cores of the even-winding and residue theorems are machine-checked in Lean~$4$. The
outstanding difficulty is isolated as a bounded-retreat / bounded-active-core lemma
together with a single-tour chronology theorem; these, and not the arithmetic
constraints proved here, are what a full solution of Conjecture~\ref{conj:main} must
still supply.

\appendix
%=======================================================================
\section{Worked example: the standard highway}\label{app:example}
%=======================================================================

The blank orbit first enters its repeating trace at step $9977$. Normalised so that
the ant faces north at the origin and the trace begins with $\turnR$ at a global
minimum of the black-count walk, the period-$104$ word (with $\turnR=$R, $\turnL=$L)
is
\begin{center}\ttfamily\small
RRRRLLRLLRRRRLLRRRRLLRLRRRRLRLLLLRRRRLRRLRRRRLLLLRLRRRRLRRRRLLLLRLRRRRLRLLRRLLLLRRLLRRRRLLRRLRLLRLLRLRLL
\end{center}
It has $\#\turnR-\#\turnL=12$ and normalised drift $(2,-2)$, in agreement with
Corollary~\ref{cor:growth4} ($\growth=12\in4\Z_{>0}$) and Corollary~\ref{cor:P16}
($12\ge 2\max(2,2)=4$). Applying Theorem~\ref{thm:criterion} to this word produces the
$13$-cell black seed
\[
\begin{aligned}
&(-6,0),\,(-5,0),\,(-5,1),\,(-4,0),\,(-4,2),\,(-3,-1),\,(-3,2),\\
&(-2,-1),\,(-2,1),\,(-1,1),\,(0,0),\,(0,1),\,(0,2),
\end{aligned}
\]
which, with a north-facing ant at the origin, produces the $104$-step highway
immediately; the leading corridor is white and the gateway pattern recurs translated
by the drift each period.

%=======================================================================
\section{Lean 4 formalization}\label{app:lean}
%=======================================================================

A dependency-light Lean~$4$ project (pinned to Lean $4.32.0$, standard library only)
machine-checks the exact transition kernel and the algebraic cores of the theorems
above. Every audited theorem depends only on the standard axioms
\texttt{propext}, \texttt{Classical.choice}, and \texttt{Quot.sound}; there are no
\texttt{sorry}s, no custom axioms, and no \texttt{native\_decide}. The checked
statements include the finite-support step kernel and its toggle lemmas, the
heading/black-count reset of Lemma~\ref{lem:invariant}, the four-corner admissibility
and mod-four fiber arithmetic of Theorem~\ref{thm:residue} (including
``one residue identity per class $\Rightarrow\growth\ge2n$''), the checker-cycle
monodromy $2\sum\alpha$ of Lemma~\ref{lem:diag}, translated-widget cancellation, and
the collision-chain endpoint parity of Theorem~\ref{thm:collision}.

The finite parity telescope \eqref{eq:telescope}---the combinatorial heart of the
even-winding theorem---is checked in the following self-contained form
(\texttt{evenWindingParityCore}); \texttt{xorAll} is the $\F_2$ sum of a list,
\texttt{incidences} takes successive XORs, and the two \texttt{false} endpoints are
the vanishing exterior half-flows:

\begin{lstlisting}[style=lean]
def bxor : Bool -> Bool -> Bool
  | false, b     => b
  | true,  false => true
  | true,  true  => false

def xorAll : List Bool -> Bool
  | []      => false
  | b :: bs => bxor b (xorAll bs)

def incidences : List Bool -> List Bool
  | []             => []
  | [_]            => []
  | a :: b :: rest => bxor a b :: incidences (b :: rest)

-- XOR of all adjacent incidences of a list bracketed by two `false`
-- endpoints is `false`: the seam-crossing count is even.
theorem evenWindingParityCore (interior : List Bool) :
    xorAll (incidences (false :: interior ++ [false])) = false := by
  rw [xorAll_incidences, lastFrom_append_singleton]; rfl
\end{lstlisting}

The formalization isolates, as an explicit interface, the endpoint normal form and
lifted potential geometry of a genuine translator, so Theorem~\ref{thm:residue} is
certified conditional on a structured trace certificate rather than end to end;
neither it nor Conjecture~\ref{conj:main} is asserted as an unconditional Lean
theorem. The project and a build script accompany the paper \cite{Repo}.

%=======================================================================
\section{Computational methodology and reproducibility}\label{app:repro}
%=======================================================================

The periodic search space is the tree of legal prefixes. By
Proposition~\ref{prop:alt} it suffices to search words beginning with $\turnR$, and by
Corollary~\ref{cor:growth4} only positive-growth heading-reset words. To distribute a
period-$N$ search, the half-open interval of length-$k$ prefix ranks $[0,2^{k-1})$ is
partitioned into disjoint worker intervals; each worker performs a depth-first
enumeration of its subtree, applying the exact criterion of
Theorem~\ref{thm:criterion} at every leaf. A complete exclusion requires that every
rank be covered exactly once, that every shard report completion with no node cap, and
that the aggregate node, leaf, and prune counters equal the sums of the per-rank
counters.

Two independent engines are run for the period $42$--$48$ certifications: a reference
implementation applying the criterion directly to every leaf, and a residue-aware
implementation that first rejects leaves violating Theorem~\ref{thm:residue} and
applies the criterion only to survivors. Agreement of all per-rank counters between
the two engines is the cross-audit. The residue identity itself was validated by an
independent enumeration of all $524{,}286$ turn words through length $18$: all
$172{,}092$ heading-reset nonzero-drift words satisfy \eqref{eq:res-a}, as do the
standard word, its powers, and its dihedral transforms. The engines' source, the
per-shard result records with SHA-256 manifests, and the exact commands to regenerate
them are provided with the paper \cite{Repo}.

%=======================================================================
\begin{thebibliography}{99}

\bibitem{BunimovichTroubetzkoy1992}
L.~A. Bunimovich and S.~E. Troubetzkoy.
\newblock Recurrence properties of {L}orentz lattice gas cellular automata.
\newblock \emph{J. Stat. Phys.} \textbf{67} (1992), no.~1--2, 289--302.

\bibitem{Etse2026confined}
K.~Etse.
\newblock How long can the escaping ant be confined?
\newblock arXiv:2606.26677 (2026).

\bibitem{GajardoLutfallaRao2024}
A.~Gajardo, V.~H. Lutfalla, and M.~Rao.
\newblock Ants on the highway.
\newblock arXiv:2409.10124 (2024); \emph{Natural Computing} \textbf{24} (2025), 497--509.

\bibitem{GajardoMoreiraGoles2002}
A.~Gajardo, A.~Moreira, and E.~Goles.
\newblock Complexity of {L}angton's ant.
\newblock \emph{Discrete Appl. Math.} \textbf{117} (2002), no.~1--3, 41--50.
\newblock arXiv:nlin/0306022.

\bibitem{Gajardo2003symbolic}
A.~Gajardo.
\newblock A symbolic projection of {L}angton's ant.
\newblock \emph{Discrete Math. Theor. Comput. Sci.} (2003).
\newblock doi:10.46298/dmtcs.2312.

\bibitem{GaleProppSutnerTroubetzkoy1995}
D.~Gale, J.~Propp, S.~Sutherland, and S.~Troubetzkoy.
\newblock Further travels with my ant.
\newblock \emph{Math. Intelligencer} \textbf{17} (1995), no.~3, 48--56.

\bibitem{Lutfalla2025sideways}
V.~H. Lutfalla.
\newblock Sideways on the highways.
\newblock arXiv:2505.05426 (2025).

\bibitem{Traldi2012transition}
L.~Traldi.
\newblock The transition matroid of a $4$-regular graph: an introduction.
\newblock arXiv:1204.0482 (2012).

\bibitem{Traldi2016binary}
L.~Traldi.
\newblock Binary nullity, {E}uler circuits and interlacement graphs.
\newblock arXiv:1607.04233 (2016).

\bibitem{Repo}
A.~Jillhewar.
\newblock Langton's ant: structural constraints, exact searches, and a Lean
formalization (supplementary code, data, and Lean project).
\newblock GitHub repository, 2026.

\end{thebibliography}

\end{document}
